% Part: incompleteness
% Chapter: representability-in-q
% Section: sigma1-completeness

\documentclass[../../../include/open-logic-section]{subfiles}

\begin{document}

\olfileid{inc}{inp}{s1c}
\olsection{\texorpdfstring{$\Sigma_1$}{Sigma-1} 完备性}

尽管 $\Th{Q}$ 及其一致、可公理化的扩张是不完全的，但我们已经看到 $\Th{Q}$ 确实能证明许多关于数词的基本事实。事实上，这可以相当大程度地推广。为了理解在~$\Th{Q}$ 中能够证明什么，我们引入 $\Delta_0$、$\Sigma_1$ 和 $\Pi_1$ 公式的概念。粗略地说，一个 $\Sigma_1$ 公式是形如 $\lexists[x][!B(x)]$ 的公式，其中 $!B$ 仅用命题联结词和有界量词构造而成。我们将证明，如果 $!A$ 是一个在 $\Struct{N}$ 中为真的 $\Sigma_1$ 语句，那么 $\Th{Q} \Proves !A$（\olref{thm:sigma1-completeness}）。

\begin{defn}
\ollabel{defn:bd-quant}
一个\emph{有界存在公式}是形如 $\lexists[x][(x < t \land !A(x))]$ 的公式，其中 $t$ 是任意项，我们通常将其记作 $\bexists{x < t}{!A(x)}$。
%
一个\emph{有界全称公式}是形如 $\lforall[x][(x < t \lif !A(x))]$ 的公式，其中 $t$ 是任意项，我们通常将其记作 $\bforall{x < t}{!A(x)}$。
\end{defn}

\begin{defn}
\ollabel{defn:delta0-sigma1-pi1-frm}
如果公式 $!B$ 是仅由原子公式经命题联结词和有界量词构造而成的，则称其为 $\Delta_0$ 的。
%
如果 $!A \ident \lexists[x][!B(x)]$，其中 $!B$ 是 $\Delta_0$ 的，则公式 $!A$ 是 $\Sigma_1$ 的。
%
如果 $!A \ident \lforall[x][!B(x)]$，其中 $!B$ 是 $\Delta_0$ 的，则公式 $!A$ 是 $\Pi_1$ 的。
\end{defn}

\begin{lem}
\ollabel{lem:q-proves-clterm-id} 设 $t$ 是一个满足 $\Value{t}{N} = n$ 的闭项。则 $\Th{Q} \Proves \eq[t][\num n]$。
\end{lem}

\begin{proof}
我们对 $t$ 的复杂度进行归纳来证明此结论。对于归纳基础，$\Value{\Obj 0}{N} = 0$，并且 $\Th{Q} \Proves \eq[\Obj 0][\num 0]$，因为 $\num 0 \ident \Obj 0$。
%
对于归纳步骤，设 $t_1$ 和 $t_2$ 是满足 $\Value{t_1}{N} = n_1$、$\Value{t_2}{N} = n_2$、$\Th{Q} \Proves \eq[t_1][\num n_1]$ 和 $\Th{Q} \Proves \eq[t_2][\num n_2]$ 的项。

那么 $\Value{(t_1')}{N} = n_1 + 1$，并且对归纳假设和公式 $\eq[\num{n_1}'][\num{n_1}']$ 应用等词的一阶规则，可得 $\Th{Q} \Proves
\eq[t_1'][{\num n_1}']$，于是根据数词的定义，我们得到 $\Th{Q} \Proves \eq[t_1'][\num{n_1 + 1}]$。

对于加法，我们有
\[
      \Value{(t_1 + t_2)}{N}
    = \Value{t_1}{N} + \Value{t_2}{N}
    = n_1 + n_2.
\]
根据归纳假设和等词的规则，$\Th{Q} \Proves \eq[t_1 + t_2][\num{n_1} + t_2]$，然后再次应用等词规则可得 $\Th{Q} \Proves \eq[t_1 + t_2][\num{n_1} + \num{n_2}]$。
根据 \olref[inc][req][bre]{lem:q-proves-add}，$\Th{Q} \Proves \eq[\num{n_1} + \num{n_2}][\num{n_1 + n_2}]$，所以 $\Th{Q} \Proves \eq[t_1 + t_2][\num{n_1 + n_2}]$。

类似的推理也适用于~$\times$，只需使用 \olref[inc][req][bre]{lem:q-proves-mult}。
%
由于这穷尽了算术中的所有闭项，因此对于所有满足 $\Value{t}{N} = n$ 的闭项~$t$，我们都有 $\Th{Q} \Proves \eq[t][\num n]$。
\end{proof}

\begin{prob}
详细证明 \olref{lem:q-proves-clterm-id} 中涉及~$\times$ 的部分。
\end{prob}

\begin{lem}
\ollabel{lem:atomic-completeness}
设 $t_1$ 和 $t_2$ 是闭项。则
\begin{enumerate}
\item 若 $\Value{t_1}{N} = \Value{t_2}{N}$，则 $\Th{Q} \Proves \eq[t_1][t_2]$。
\item 若 $\Value{t_1}{N} \neq \Value{t_2}{N}$，则 $\Th{Q} \Proves \eq/[t_1][t_2]$。
\item 若 $\Value{t_1}{N} < \Value{t_2}{N}$，则 $\Th{Q} \Proves t_1 < t_2$。
\item 若 $\Value{t_2}{N} \leq \Value{t_1}{N}$，则 $\Th{Q} \Proves \lnot(t_1 < t_2)$。
\end{enumerate}
\end{lem}

\begin{proof}
给定项 $t_1$ 和 $t_2$，我们取定 $n = \Value{t_1}{N}$ 和 $m = \Value{t_2}{N}$。

假设 $!A \ident t_1 = t_2$。根据 \olref{lem:q-proves-clterm-id}，$\Th{Q} \Proves \eq[t_1][\num n]$ 且 $\Th{Q} \Proves \eq[t_2][\num n]$。
如果 $n = m$，则 $\Th{Q} \Proves \eq[\num n][\num m]$，从而根据等词的传递性有 $\Th{Q} \Proves \eq[t_1][t_2]$。
如果 $n \neq m$，则 $\Th{Q} \Proves \eq/[\num n][\num m]$，并再次根据等词的传递性，有 $\Th{Q} \Proves \eq/[t_1][t_2]$。

现在令 $!A \ident t_1 < t_2$。对于两种情况，我们都依赖公理~$!Q_8$，该公理断言对所有的 $x,y$ 都有 $x < y \liff \lexists[z][\eq[z' + x][y]]$。

假设 $\Sat{N}{t_1 < t_2}$。则存在某个 $k \in \Nat$ 使得 $n + k + 1 = m$。根据 \olref{lem:q-proves-clterm-id}，$\Th{Q} \Proves \eq[t_1][\num n]$ 且 $\Th{Q} \Proves \eq[t_2][\num m]$，并且根据本引理的第一部分，$\Th{Q} \Proves \eq[\num n + {\num k}'][\num m]$。
根据等词的传递性可得 $\Th{Q} \Proves \eq[{\num k}' + t_1][t_2]$，所以 $\Th{Q} \Proves \lexists[z][\eq[z' + t_1][t_2]]$。
根据~$!Q_8$ 的从右到左方向，$\Th{Q} \Proves t_1 < t_2$。

假设 $\Sat/{N}{t_1 < t_2}$，即 $m \leq n$。
%
我们在~$\Th{Q}$ 中工作，并假设 $t_1 < t_2$。根据~$!Q_8$ 的从左到右方向，存在某个~$z$ 使得 $\eq[z' + t_1][t_2]$。
由于 $\Th{Q} \Proves \eq[t_1][\num n]$ 且 $\Th{Q} \Proves \eq[t_2][\num m]$，故 $\eq[z' + \num n][\num m]$。
%
利用~$!Q_5$ 对~$m$ 进行外部归纳，可得 $\eq[z' + \num{n - m}][\Obj 0]$。
如果 $m = n$，则 $\eq/[z'][\Obj 0]$，从而由~$!Q_3$ 得出矛盾。
如果 $m < n$，则再次由~$!Q_5$ 得到 $\eq[(z' + \num{n - m - 1})'][\Obj 0]$，从而由~$!Q_3$ 得出矛盾。
因此 $\Th{Q} \Proves \lnot(t_1 < t_2)$。
\end{proof}

\begin{lem}
\ollabel{lem:bounded-quant-equiv}
设 $!A$ 是一个公式，$t$ 是一个闭项，且 $k=\Value{t}{N}$。则
\begin{enumerate}
\item $\Th{Q} \Proves \bforall{x<t}{!A(x)}$ 当且仅当 $\Th{Q} \Proves
    !A(\num 0) \land \dots \land !A(\num{k-1})$。
\item $\Th{Q} \Proves \bexists{x<t}{!A(x)}$ 当且仅当 $\Th{Q} \Proves
    !A(\num 0) \lor \dots \lor !A(\num{k-1})$。
\end{enumerate}
\end{lem}

\begin{proof}
    我们证明有界全称量词的情形。
    如果 $\Value{t}{N} = 0$，那么等价式的左边在~$\Th{Q}$ 中是可证的，因为根据 \olref[inc][req][min]{lem:less-zero}，不存在 $x<\num 0$。
    类似地，我们可以将空析取简单地视为 $\ltrue$，这在~$\Th{Q}$ 中也是可证的。
    %
    因此假设存在某个自然数 $k$ 使得 $\Value{t}{N} = k+1$。根据 \olref{lem:q-proves-clterm-id}，我们可以假定讨论的是形如 $\bforall{x<\num{k+1}}{!A(x)}$ 的公式。
    
    假设 $\Th{Q} \Proves \bforall{x<\num{k+1}}{!A(x)}$，
    并令 $n \leq k$。由于根据 \olref{lem:atomic-completeness} 有 $\Th{Q} \Proves \num n < \num{k+1}$，根据逻辑可知 $\Th{Q} \Proves !A(\num n)$。对每个 $n \leq k$ 应用这一事实 $k+1$ 次，我们便得到 $\Th{Q} \Proves !A(\num 0)
    \land \dots \land !A(\num k)$，此即所求。
    
    对于另一个方向，假设 $\Th{Q} \Proves
    !A(\num 0) \land \dots \land !A(\num k)$。在 $\Th{Q}$ 中工作，假设 $x < \num{k+1}$。
    根据 \olref[inc][req][min]{lem:less-nsucc}，我们有 $x = \num 0 \lor \dots \lor x = \num k$，因此根据逻辑可知~$!A(x)$，从而可得全称命题 $\bforall{x<\num{k+1}}{!A(x)}$。
    
    有界存在量词公式的等价性证明是类似的。
\end{proof}

\begin{prob}
给出 \olref{lem:bounded-quant-equiv} 中存在量词情形的详细证明。
\end{prob}

\begin{lem}
\ollabel{lem:delta0-completeness}
如果 $!A$ 是一个在 $\Struct{N}$ 中为真的 $\Delta_0$ 语句，那么 $\Th{Q} \Proves !A$。
\end{lem}

\begin{proof}
我们通过对公式复杂度归纳来证明此引理。
%
归纳基础由 \olref{lem:atomic-completeness} 给出，所以我们转向归纳步骤。为简便起见，我们根据否定符号所应用的公式的结构，将否定情形分成子情况处理。

\begin{enumerate}
\item 假设 $(!A \land !B)$ 在 $\Struct{N}$ 中为真，因此 $!A$ 和 $!B$ 在~$\Struct{N}$ 中为真。
根据归纳假设，$\Th{Q} \Proves !A$ 且 $\Th{Q} \Proves !B$，因此由逻辑可得 $\Th{Q} \Proves (!A \land !B)$。
%
\item 假设 $\lnot (!A \land !B)$ 在 $\Struct{N}$ 中为真，因此 $\lnot !A$ 或 $\lnot !B$ 在 $\Struct{N}$ 中为真。
不失一般性，假设前者成立。根据归纳假设 $\Th{Q} \Proves \lnot !A$，因此由逻辑可得 $\Th{Q} \Proves \lnot (!A \land !B)$。
%
\item 假设 $(!A \lor !B)$ 在 $\Struct{N}$ 中为真，因此 $!A$ 或 $!B$ 在 $\Struct{N}$ 中为真。不失一般性，假设前者成立。根据归纳假设 $\Th{Q} \Proves !A$，因此由逻辑可得 $\Th{Q} \Proves (!A \lor !B)$。
%
\item 假设 $\lnot(!A \lor !B)$ 在 $\Struct{N}$ 中为真，因此在 $\Struct{N}$ 中 $\lnot !A$ 和 $\lnot !B$ 均为真。
那么根据归纳假设，$\Th{Q} \Proves \lnot !A$ 且 $\Th{Q} \Proves \lnot !B$。因此由逻辑可得 $\Th{Q} \Proves \lnot(!A \lor !B)$。
%
\item 假设 $\bforall{x<t}{!A(x)}$ 在~$\Struct{N}$ 中为真，其中 $t$ 是闭项且 $k=\Value{t}{N}$。根据归纳假设和逻辑，如果对每个 $n < \Value{t}{N}$，$!A(\num n)$ 都在~$\Struct{N}$ 中为真，那么 $\Th{Q} \Proves
!A(\num 0) \land \dots \land !A(\num{k-1})$。
根据 \olref{lem:bounded-quant-equiv} 可得 $\Th{Q} \Proves \bforall{x<t}{!A(x)}$。
%
\item 有界存在量词的情形，即形如 $\bexists{x < t}{!A(x)}$ 的语句，与有界全称量词的情形类似。
%
\item 假设 $\lnot \bforall{x<t}{!A(x)}$ 在~$\Struct{N}$ 中为真，其中 $t$ 是闭项。此语句等价于语句 $\bexists{x<t}{\lnot !A(x)}$，且该等价关系在~$\Th{Q}$ 中可导，因此我们可以应用关于有界存在量词的推理。
%
\item 类似地，假设 $\lnot \bexists{x<t}!A(x)$ 在 $\Struct{N}$ 中为真，其中 $t$ 是闭项。此语句在 $\Th{Q}$ 中等价于 $\bforall{x<t}{\lnot!A(x)}$，因此我们可以应用关于有界全称量词的推理。
%
\item 最后，假设 $\lnot !A$ 在 $\Struct{N}$ 中为真。
剩下的情形只有 $!A$ 是原子公式，以及 $\lnot !A \ident \lnot\lnot !B$（对某个 $\Delta_0$ 语句 $!B$）两种。如果 $!A$ 是原子公式，那么根据 \olref{lem:atomic-completeness}，$\Th{Q} \Proves \lnot !A$。如果 $\lnot !A \ident \lnot\lnot !B$，那么根据逻辑，$\lnot\lnot !B$ 在~$\Th{Q}$ 中可证等价于~$!B$；而由于 $\lnot !A$ 在~$\Struct{N}$ 中为真，$!B$ 也在~$\Struct{N}$ 中为真。根据归纳假设，我们因此有 $\Th{Q} \Proves \lnot !A$。
\end{enumerate}
\end{proof}

\begin{prob}
给出 \olref{lem:delta0-completeness} 中存在量词情形的详细证明。
\end{prob}

\begin{thm}
\ollabel{thm:sigma1-completeness}
如果 $!A$ 是一个在~$\Struct{N}$ 中为真的 $\Sigma_1$ 语句，那么 $\Th{Q} \Proves !A$。
\end{thm}

\begin{proof}
如果 $\lexists{x}!A(x)$ 是一个在~$\Struct{N}$ 中为真的 $\Sigma_1$ 语句，
则存在一个自然数~$n$ 和一个变元指派~$s$，使得 $s(x) = n$ 并且
$\Sat{N}{!A(x)}[s]$。
根据关于满足关系的标准性质可得，
$\Sat{N}{!A(\num n)}$。
但 $!A(\num n)$ 是一个 $\Delta_0$ 公式，
因此由 \olref{lem:delta0-completeness}，我们有
$\Th{Q} \Proves !A(\num n)$，
从而由逻辑我们也有
$\Th{Q} \Proves \lexists[x][!A(x)]$。
\end{proof}

\end{document}